%--- iliad ---
% separateSolutions: true
% title: Condensation
% summary: >-
%   Organizing information across a family of observations. Define latent
%   models, examine reconstruction and conditional independence, and prove
%   an exact and quantitative correspondence theorem.
%--- end ---
\documentclass[11pt]{article}
\IfFileExists{iliad.sty}{\usepackage{iliad}}{\usepackage{../iliad}}
\geometry{a4paper,margin=25mm}
\usepackage{microtype}
\AtBeginEnvironment{example}{\crefalias{theorem}{example}}
\AtBeginEnvironment{definition}{\crefalias{theorem}{definition}}
\AtBeginEnvironment{lemma}{\crefalias{theorem}{lemma}}
\AtBeginEnvironment{corollary}{\crefalias{theorem}{corollary}}
\AtBeginEnvironment{proposition}{\crefalias{theorem}{proposition}}
\newcommand{\PI}{\mathcal{P}^{+}(I)}
\newcommand{\Up}{\mathord{\supseteq}}
\title{Condensation}
\author{Satya Benson}
\date{21 September 2026}
\begin{document}
\maketitle

\section{Prerequisites}
\label{sec:prerequisites}
We use probability on finite sets, conditional probability, and elementary
proofs. All random variables in this sheet have finite ranges. All logarithms
have base two.

\begin{learningoutcomes}
\begin{itemize}
\item Construct a latent model and distinguish LVM from reconstruction.
\item Explain what the Markov condition excludes, using a concrete example.
\item Prove correspondence between collections of latents, first exactly and
then with explicit error terms.
\item Distinguish the correspondence theorem from claims about learning,
computationally finding a translation, or choosing observables.
\end{itemize}
\end{learningoutcomes}

\subsection*{How to use this sheet}
Read \cref{sec:questions,sec:conditions} after the first lecture, then work on
exercises E1--E3. After the second lecture, read
\cref{sec:correspondence} and complete E4--E8. These eight exercises form the
core route. E9 and E10 are optional extensions.

\subsection*{Information theory reference}
For a finite-valued variable $U$,
\[
 H(U)=-\sum_u p(u)\log p(u),\qquad
 H(U\mid V)=\sum_v p(v)H(U\mid V=v),
\]
with $0\log 0=0$. Conditional entropy is nonnegative. It is zero precisely when
$U$ is a function of $V$ almost surely: values on probability-zero events do not
matter. The chain rule and conditioning inequality are
\[
 H(U,V\mid W)=H(U\mid W)+H(V\mid U,W),\qquad
 H(U\mid V,W)\le H(U\mid W).
\]
Conditional mutual information is
\[
 I(U;V\mid W)=H(U\mid W)-H(U\mid V,W)\ge0.
\]
It is symmetric in $U,V$ and is zero precisely when $U,V$ are conditionally
independent given $W$. Omitting $W$ gives mutual information $I(U;V)$.
Entropy of a tuple means joint entropy, not the sum of its marginal entropies.
An empty tuple is constant and has zero entropy.
For probability distributions $p,q$ on a finite set,
\[
 D_{\mathrm{KL}}(p\Vert q)=\sum_u p(u)\log\frac{p(u)}{q(u)}\ge0.
\]
This last quantity appears only in the optional score discussion.

\section{A family of questions}
\label{sec:questions}
Fix jointly distributed observables $X_1,\ldots,X_n$. We look for latent
variables, each assigned to a set of observations, such that the latents
assigned to an observation determine it.
We will give conditions under which corresponding latent collections in two
such models determine one another, and bound the remaining conditional entropy
when the conditions hold approximately. We take the joint distributions as given.

\begin{remark}[Choosing the observables]
One can write $X_i=f_i(S)$ for an underlying variable $S$ representing sense
data or a physical state. This introduces no restriction: $S$ could simply be
the tuple of all observations. An observable distinguishes outcomes by
partitioning them into its possible answers; random-variable values label those
parts. Information quantities do not depend on bijective relabeling of answers.
The choice of which observables to include, and their joint law, remains part
of the model.
\end{remark}

\begin{example}[Overlapping observations]
\label{eg:bits}
Let $S,P,Q$ be independent fair bits and set
\[
 X_1=(S,P),\qquad X_2=(S,Q),\qquad X_3=(P,Q).
\]
Use a latent $Y_{\{1,2\}}=S$ for observations 1 and 2, a latent
$Y_{\{1,3\}}=P$ for observations 1 and 3, and a latent $Y_{\{2,3\}}=Q$
for observations 2 and 3. Each observation is determined by the two latents
assigned to it. Each of these latents can also be recovered from either
observation to which it is assigned.
\end{example}

\subsection{Addresses and collections}
Let $I=\{1,\ldots,n\}$ and write $\PI$ for its nonempty subsets.
An \emph{address} $A\in\PI$ specifies which observations receive a latent
$Y_A$. The collection of addresses is a partially ordered set under inclusion.
We use one slot for each nonempty subset, setting unused slots to constants.

\begin{definition}[Latent model]
\label{def:model}
A latent model of $X=(X_i)_{i\in I}$ consists of finite-valued variables
$Y=(Y_A)_{A\in\PI}$ jointly distributed with $X$, such that
\begin{equation}
 H\bigl(X_i\mid Y_{\{A\in\PI\,:\,i\in A\}}\bigr)=0
 \quad\text{for every }i\in I.
 \label{eq:forward}
\end{equation}
Here $Y_{\mathcal{F}}=(Y_A)_{A\in\mathcal{F}}$ for a family
$\mathcal{F}\subseteq\PI$.
\end{definition}

We call \eqref{eq:forward} the \emph{latent variable model condition}, or
\emph{LVM}. The contributing latents, taken together, determine the
observation. This is not a claim that an agent knows their realized values. Some slots may represent residual noise.
Also, an address need not be minimal: being assigned to an observation does not
assert that a slot makes a nonredundant or causal contribution to it.

For nonempty $A\subseteq I$ and $B\subseteq I$, define two different families:
\[
 \Up A=\{C\in\PI:A\subseteq C\},\qquad
 \mathcal{J}_B=\{C\in\PI:C\cap B\ne\varnothing\}.
\]
The tuple $Y_{\Up A}$ is the \emph{tower above $A$}. The tuple
$Y_{\mathcal{J}_B}$ contains all latents contributing to at least one observation
in $B$, and therefore determines $X_B=(X_i)_{i\in B}$.
In particular $\Up\{i\}=\mathcal{J}_{\{i\}}$ is the family of addresses
contributing to $X_i$, so \eqref{eq:forward} reads $H(X_i\mid Y_{\Up\{i\}})=0$.
For a larger $A$, $\Up A$ and $\mathcal{J}_A$ generally differ.

\begin{exercise}[E1: Overlapping bits; 15 minutes]
\label{ex:bits}\important
Use \cref{eg:bits}, with all unmentioned slots constant.
\begin{enumerate}
\item Compute $H(X_1)$, $H(X_1,X_2,X_3)$, and $I(X_1;X_2)$.
\item List $\Up\{1\}$, $\Up\{2\}$, and $\Up\{1,2\}$, including
constant slots. Which tuple determines $X_1$? Which of the three is
$\Up\{1\}\cap\Up\{2\}$?
\item Give explicit functions recovering $X_1$ from its contributing latents
and recovering $Y_{\{1,2\}}$ from $X_1$.
\item Does the address structure have to be a tree? Identify two overlapping
addresses in this example neither of which contains the other.
\end{enumerate}
\end{exercise}

\section{Two conditions on a latent model}
\label{sec:conditions}
LVM is not one of the two conditions below. It is part of the definition of a
latent model, and every observable law admits one by construction: put a copy
of each $X_i$ in its singleton slot, or the entire tuple $X$ in the top slot
$Y_I$, and LVM holds either way. Satisfying it therefore says nothing about how
information is shared. The two conditions below constrain these alternatives in
different ways.

\begin{definition}[Reconstruction]
\label{def:reconstruction}
A latent model has exact reconstruction if
\[
 H(Y_A\mid X_i)=0\qquad\text{whenever }i\in A.
\]
Equivalently, $H(Y_{\Up A}\mid X_i)=0$ whenever $i\in A$: every slot
in that finite tuple can be recovered from $X_i$.
\end{definition}

The equivalence in this definition uses exact zero errors. In approximate
statements we will bound the entropy of the \emph{whole tower}; separately small
errors for its individual slots could add up.

A family $\mathcal{F}\subseteq\PI$ is \emph{upward closed} if
$A\in\mathcal{F}$ and $A\subseteq C$ imply $C\in\mathcal{F}$.
For example, $\Up A$ and $\mathcal{J}_B$ are upward closed. Taking a
collection together with every slot above it prevents us from omitting higher
variables on which it may depend.

\begin{definition}[Ordered Markov condition and defect]
\label{def:markov}
The ordered Markov condition is
\[
 I(Y_{\mathcal{F}};Y_{\mathcal{G}}\mid
 Y_{\mathcal{F}\cap\mathcal{G}})=0
\]
for all upward-closed families $\mathcal{F},\mathcal{G}\subseteq\PI$.
Its quantitative defect is
\[
 \eta_Y=\max_{\mathcal{F},\mathcal{G}\text{ upward closed}}
 I(Y_{\mathcal{F}};Y_{\mathcal{G}}\mid Y_{\mathcal{F}\cap\mathcal{G}}).
\]
\end{definition}

The conditioning family $\mathcal{F}\cap\mathcal{G}$ is an intersection of
\emph{address families}, not a quantity inferred to be common by inspecting
their values. The condition says that the latents at those shared addresses
account for the dependence between the two collections. It does not require
every latent to be independent of every other latent. In the two-observable
case, its only potentially nonzero requirement is
\[
 I(Y_{\{1\}};Y_{\{2\}}\mid Y_{\{1,2\}})=0.
\]
This is a graphical-model condition for the inclusion order. The equivalent
local version says that each slot is conditionally independent of all
incomparable slots given its strict supersets.

\begin{definition}[Perfect condensation]
\label{def:perfect}
A latent model is a perfect condensation if it has exact reconstruction
and satisfies the ordered Markov condition.
\end{definition}
This is equivalent to Eisenstat's definition by optimal conditioned scores
(Theorem 5.10 of the paper). \Cref{sec:scores} explains that connection.
The overlapping-bit model in \cref{eg:bits} is a perfect condensation.

\begin{exercise}[E2: Two failed models; 20 minutes]
\label{ex:failed}\important
Keep the observables of \cref{eg:bits}. Consider two alternative models.
\begin{enumerate}
\item Set $Y_{\{i\}}=X_i$ for $i=1,2,3$ and every other slot constant.
Check LVM and reconstruction. Take
$\mathcal{F}=\Up\{1\}$ and $\mathcal{G}=\Up\{2\}$.
Calculate the conditional mutual information in \cref{def:markov}.
Which information is duplicated across singleton slots?
\item Instead set $Y_{\{1,2,3\}}=(S,P,Q)$ and every other slot constant.
Check LVM and the Markov condition. Compute
$H(Y_{\{1,2,3\}}\mid X_1)$. Which condition fails?
\item Explain why LVM, reconstruction, and the
Markov condition have different roles.
\end{enumerate}
\end{exercise}

\subsection{Existence is a separate question}
\label{sec:existence}
The preceding examples test particular latent models. A stronger question is
whether a given observable law admits \emph{any} perfect condensation.
Even a pair of correlated binary variables can fail to do so.

Let $I=\{1,2\}$, let $X_1$ be a fair bit, and set $X_2=X_1\mathbin{\oplus}N$,
where $N$ is independent Bernoulli$(q)$ noise and $0<q<1/2$. Here
$\oplus$ is addition modulo two, so $a\oplus b$ is $0$ when the two bits agree
and $1$ when they differ.
The joint law is
\[
\renewcommand{\arraystretch}{1.8}
\begin{array}{r|cc}
 P(x_1,x_2) & \;\;x_2=0\;\; & \;\;x_2=1\;\;\\\hline
 x_1=0 & \dfrac{1-q}{2} & \dfrac{q}{2}\\
 x_1=1 & \dfrac{q}{2} & \dfrac{1-q}{2}
\end{array}
\]
Every one of the four outcomes has positive probability, because $0<q<1/2$;
and $X_1,X_2$ are dependent, because the two rows differ.
Suppose a perfect condensation existed, and consider its top slot
$Y_{\{1,2\}}$.

With two observations there are only three addresses, so the whole model is
three slots:
\[
\begin{array}{r|ccc}
 \text{address} & \{1\} & \{1,2\} & \{2\}\\
 \text{slot} & Y_{\{1\}} & Y_{\{1,2\}} & Y_{\{2\}}
\end{array}
\]
LVM says $X_1$ is determined by $(Y_{\{1\}},Y_{\{1,2\}})$ and $X_2$ by
$(Y_{\{2\}},Y_{\{1,2\}})$. The top slot $Y_{\{1,2\}}$ is the only one assigned to both
observations, so it is the only place shared information can sit.

Reconstruction would give $Y_{\{1,2\}}=f(X_1)=g(X_2)$ almost surely. Because
all four observable pairs have positive probability, these equalities force
\[
 f(0)=g(0)=g(1)=f(1).
\]
Thus $Y_{\{1,2\}}$ is constant. The Markov condition would then require
$Y_{\{1\}}$ and $Y_{\{2\}}$ to be independent. LVM
makes each $X_i$ a function of its singleton slot and the constant $Y_{\{1,2\}}$.
Functions of independent variables are independent, contradicting the chosen
observable law. No perfect condensation exists for this pair.

There is a quantitative obstruction if we retain \emph{exact} reconstruction. The same argument makes $Y_{\{1,2\}}$ constant. Data processing then
gives
\[
 \eta_Y=I(Y_{\{1\}};Y_{\{2\}})
 \ge I(X_1;X_2)=1-h_2(q)>0,
 \qquad h_2(q)=-q\log q-(1-q)\log(1-q).
\]
Dependence between the observables must remain as a Markov defect when their
separately recoverable top slot is constant.

That bound pins one end of a trade-off: exact reconstruction forces a Markov
defect of at least $1-h_2(q)$, and tolerating some reconstruction error buys
part of it back. When no perfect condensation exists, which is the usual
situation, the latent models of a given observable law trace out a Pareto
frontier between reconstruction defect and Markov defect. The useful question
is then whether some point on that frontier has both low \emph{enough} for the
correspondence bound of \cref{sec:correspondence} to say something.

\begin{remark}
One could instead ask for the model minimizing the conditioned score of
\cref{sec:scores}. The score is defined per observation block, however, and
there is not always a single model minimizing it everywhere.
\end{remark}

\begin{pdfonly}
\pagebreak
\end{pdfonly}

\begin{exercise}[E3: Collections, not individual slots; 15 minutes]
\label{ex:duplicates}\important
For this exercise let $I=\{1,2\}$, take independent fair bits $C,R,T$, and set
$X_1=(C,R)$, $X_2=(C,T)$. Compare
\[
\begin{array}{lll}
 Y_{\{1,2\}}=C,&Y_{\{1\}}=R,&Y_{\{2\}}=T,\\
 Z_{\{1,2\}}=C,&Z_{\{1\}}=(C,R),&Z_{\{2\}}=(C,T).
\end{array}
\]
\begin{enumerate}
\item Check that both models are perfect condensations.
\item Compute $H(Z_{\{1\}}\mid Y_{\{1\}})$.
\item Show that $Y_{\Up\{1\}}$ and $Z_{\Up\{1\}}$ determine one
another. Explain why correspondence between these tuples is a different
claim from correspondence between the singleton slots.
\end{enumerate}
\end{exercise}

\section{The correspondence theorem}
\label{sec:correspondence}
Take two latent models of the same observable tuple $X$. To compare them,
put $(X,Y,Z)$ on one probability space while retaining both specified laws
$(X,Y)$ and $(X,Z)$. For finite variables this is always possible: for example,
use $p(x,y,z)=p(x)p(y\mid x)p(z\mid x)$. The theorem holds for any such
coupling.

\begin{theorem}[Correspondence]
\label{thm:correspondence}
For two latent models $Y,Z$ of $X$ and any nonempty $A\subseteq I$, let
$k=|A|$. Then
\begin{equation}
 H(Y_{\Up A}\mid Z_{\Up A})
 \le \sum_{i\in A}H(Y_{\Up A}\mid X_i)+(k-1)\eta_Z.
 \label{eq:correspondence}
\end{equation}
If each term in the sum is at most $\delta$ and $\eta_Z\le\eta$, the bound
is $k\delta+(k-1)\eta$. If both budgets are $\varepsilon$, it is
$(2k-1)\varepsilon$. Interchanging $Y,Z$ gives the reverse bound.
\end{theorem}

\begin{corollary}[Exact correspondence]
\label{cor:exact}
If $Y$ and $Z$ are perfect condensations, their towers above every nonempty
address $A$ determine one another almost surely:
\[
 H(Y_{\Up A}\mid Z_{\Up A})=H(Z_{\Up A}\mid Y_{\Up A})=0.
\]
\end{corollary}

The forward bound uses reconstruction in the sender $Y$, LVM in the receiver
$Z$, and the Markov defect of $Z$.
It does not use the Markov condition of $Y$ until we reverse the argument.
The result concerns corresponding collections. E3 shows why we cannot simply
replace those collections by individual slots.

This proof is a specialization of Eisenstat's approximate correspondence
theorem (Theorem 6.8), using the two-witness inequality (Lemma 6.4).
Its tower presentation follows the exposition listed in \cref{sec:reading}.
The exercises below give the full argument without additional machinery.

\subsection{Proof roadmap}
The theorem has three steps.
\begin{enumerate}
\item If two witnesses each determine a target and are independent given their
common part, that common part determines the target. Entropy gives a version
with errors.
\item Each observation $X_i$ is determined by its contributing $Z$-latents.
Anything recoverable from $X_i$ is therefore recoverable from that collection.
\item Intersect these contributing collections over $i\in A$. Exactly the
addresses containing all of $A$ survive. Each intersection costs at most one
Markov error term.
\end{enumerate}
The word ``witness'' below means a variable or tuple from which the target is
recoverable.

\begin{lemma}[Two witnesses]
\label{lem:witness}
For any finite-valued variables $U,M,P,Q$,
\begin{equation}
 H(U\mid M)\le H(U\mid M,P)+H(U\mid M,Q)+I(P;Q\mid M).
 \label{eq:witness}
\end{equation}
\end{lemma}

\begin{exercise}[E4: Prove the two-witness lemma; 20 minutes]
\label{ex:witness}\important
\begin{enumerate}
\item Prove \cref{lem:witness} by expanding conditional entropies using the
chain rule. Find an expression for the difference between the right and left
sides that is a sum of nonnegative quantities.
\item Deduce the exact statement when the three terms on the right are zero.
\item Explain the three error terms. Which measure failures of reconstruction,
and which measures dependence that remains after revealing $M$?
\end{enumerate}
\end{exercise}
\begin{hint}
Try expanding $I(P;Q\mid M)-I(P;Q\mid U,M)$, then rearranging with the
chain rule, which is stated in the
\href{\#information-theory-reference}{information theory reference}.
The desired difference contains a conditional mutual information
and a conditional entropy.
\end{hint}

\begin{exercise}[E5: The dependence term is needed; 5 minutes]
\label{ex:dependence}
Let $U=P=Q$ be a fair bit and let $M$ be constant. Evaluate every term in
\eqref{eq:witness}. What would go wrong if the mutual-information term were
omitted? How does this example resemble the singleton-copy model in E2?
\end{exercise}

\begin{exercise}[E6: Transfer through an observable; 10 minutes]
\label{ex:transfer}\important
Let $U$ be any finite-valued target and suppose $X_i$ is determined by
$Z_{\Up\{i\}}$. Prove
\[
 H(U\mid Z_{\Up\{i\}})\le H(U\mid X_i).
\]
State the equality obtained by adding $X_i$ to the conditioning tuple, then
apply the conditioning inequality. Does this argument assume that $U$ and
$Z_{\Up\{i\}}$ are conditionally independent given $X_i$?
\end{exercise}

\begin{exercise}[E7: Intersect the address families; 15 minutes]
\label{ex:intersections}\important
\begin{enumerate}
\item For $I=\{1,2,3\}$, list $\Up\{1\}$, $\Up\{2\}$, and
$\Up\{3\}$. Compute $\Up\{1\}\cap\Up\{2\}$ and then its
intersection with $\Up\{3\}$.
\item Show that an intersection of upward-closed families is upward closed.
\item For general nonempty $A\subseteq I$, prove
$\bigcap_{i\in A}\Up\{i\}=\Up A$ by checking membership of an address
$C$ on both sides.
\item If $\mathcal{F},\mathcal{G}$ are upward closed, apply the two-witness
lemma with $M=Z_{\mathcal{F}\cap\mathcal{G}}$,
$P=Z_{\mathcal{F}}$, and $Q=Z_{\mathcal{G}}$.
Why can the extra conditioning on $M$ be dropped from each reconstruction
term? Write the resulting bound using $\eta_Z$.
\end{enumerate}
\end{exercise}

\begin{pdfonly}
\pagebreak
\end{pdfonly}

\begin{exercise}[E8: Complete the correspondence proof; 25 minutes]
\label{ex:induction}\important
Fix $A=\{i_1,\ldots,i_k\}$, set
$\mathcal{G}_j=\bigcap_{\ell=1}^{j}\Up\{i_\ell\}$.
\begin{enumerate}
\item Use E6 to prove the bound at $j=1$.
\item Use E7 to show, for $j\ge2$,
\[
 H(Y_{\Up A}\mid Z_{\mathcal{G}_j})\le
 H(Y_{\Up A}\mid Z_{\mathcal{G}_{j-1}})+H(Y_{\Up A}\mid Z_{\Up\{i_j\}})+\eta_Z.
\]
\item Prove by induction that
\[
 H(Y_{\Up A}\mid Z_{\mathcal{G}_j})\le
 \sum_{\ell=1}^{j}H(Y_{\Up A}\mid X_{i_\ell})+(j-1)\eta_Z.
\]
Evaluate this at $j=k$ to prove \cref{thm:correspondence}. Check the
case $k=1$ explicitly.
\item Deduce \cref{cor:exact}, including the reverse direction. Name the
assumption used at each step of the proof.
\item Suppose instead that $H(Y_C\mid X_i)\le\varepsilon$ separately for
all relevant slots $C$. Why does this not by itself supply the tower error
budget $H(Y_{\Up A}\mid X_i)\le\varepsilon$?
\end{enumerate}
\end{exercise}

\section{Optional: several observations as one witness}
\label{sec:blocks}
A shared parameter need not be recoverable from a single observation. For
example, take $n\ge2$, let $C$ be a fair bit, and set
$X_i=C\mathbin{\oplus}N_i$, where the
$N_i$ are independent Bernoulli$(p)$ noises with $0<p<1/2$, independent of
$C$. The model $Y_I=C$, $Y_{\{i\}}=X_i$ satisfies LVM and the
Markov condition exactly, but $H(C\mid X_i)=h_2(p)>0$, where
$h_2(p)=-p\log p-(1-p)\log(1-p)$.
Larger observation blocks can reveal more about $C$. To use them in the proof,
replace $\Up\{i\}$ by $\mathcal{J}_B$.

\begin{exercise}[E9: Block witnesses; optional, 25 minutes]
\label{ex:blocks}
\begin{enumerate}
\item Let $I=\{1,2,3\}$ and suppose a target $U$ has conditional entropy at
most $\delta$ given each of $X_{\{1,2\}}$, $X_{\{1,3\}}$, and
$X_{\{2,3\}}$. Which addresses survive the intersection
$\mathcal{J}_{\{1,2\}}\cap\mathcal{J}_{\{1,3\}}\cap\mathcal{J}_{\{2,3\}}$?
Prove a bound on the entropy of $U$ given the surviving $Z$-slots.
\item More generally, let $\mathcal{W}$ be a nonempty finite family of nonempty
observation blocks. Adapt E8 to show
\[
 H(U\mid Z_{\mathcal{R}})\le
 \sum_{B\in\mathcal{W}}H(U\mid X_B)+(|\mathcal{W}|-1)\eta_Z,
 \qquad \mathcal{R}=\bigcap_{B\in\mathcal{W}}\mathcal{J}_B.
\]
\item Fix an integer $0\le r<|A|$ and take every $(r+1)$-element subset of $A$ as a
witness. Show that $C\in\mathcal{R}$ exactly when $|A\setminus C|\le r$.
What does this give for $r=0$? Why is the receiver generally larger for $r>0$?
\end{enumerate}
\end{exercise}

Eisenstat's later almost-perfect-condensation formulation bounds the tower
reconstruction error whenever $|A\cap B|>r$, the Markov defect, and the number
of active slots. The threshold allows small slots to account for local noise
without requiring every such slot to be recovered from every relevant block.
It also changes which receiver the proof isolates, as E9 shows. The exact-tower
bound with singleton witnesses does not follow from an arbitrary positive
threshold $r$.

\section{Optional: the conditioned score}
\label{sec:scores}
The conditions also have a description-length interpretation. For an
observation block $B$, compare
\[
 \sigma(B)=\sum_{A\in\mathcal{J}_B}H(Y_A),\qquad
 \chi(B)=\sum_{A\in\mathcal{J}_B}H(Y_A\mid Y_{\Up A\setminus\{A\}}).
\]
The simple score $\sigma$ charges each active slot separately. The conditioned
score $\chi$ charges its additional information given strict supersets.
Those supersets are also active for $B$. Encoding in reverse inclusion order
makes this conditioning available. These are ideal expected description-length
quantities; they do not measure computational retrieval time.

Both scores are at least $H(X_B)$. More precisely, put
$V_B=\chi(B)-H(Y_{\mathcal{J}_B})$. Then
\begin{equation}
 \chi(B)-H(X_B)=V_B+H(Y_{\mathcal{J}_B}\mid X_B),\qquad V_B\ge0.
 \label{eq:score}
\end{equation}
The second term charges latent information not recoverable from $X_B$.
The first charges dependence not accounted for by the inclusion hierarchy.
It vanishes under the ordered Markov condition.

For a graphical-model expression, define
\[
 q_B(y)=\prod_{A\in\mathcal{J}_B}
 p(y_A\mid y_{\Up A\setminus\{A\}}).
\]
This is a normalized distribution: sample larger addresses before smaller
ones. Conditional kernels on probability-zero parent configurations can be
chosen arbitrarily. Direct expansion gives
$V_B=D_{\mathrm{KL}}(p_{Y_{\mathcal{J}_B}}\Vert q_B)$.
Thus \eqref{eq:score} separates two nonnegative defects.

\begin{exercise}[E10: The two charges; optional, 25 minutes]
\label{ex:scores}
\begin{enumerate}
\item Order $\mathcal{J}_B$ so strict supersets appear first. Use the chain rule
and conditioning inequality to prove $\chi(B)\ge H(Y_{\mathcal{J}_B})$.
\item Use LVM to prove
$H(Y_{\mathcal{J}_B})-H(X_B)=H(Y_{\mathcal{J}_B}\mid X_B)$, hence
\eqref{eq:score}.
\item For both models in E3, calculate $\sigma(\{1,2\})$ and
$\chi(\{1,2\})$. Why is copying the top variable into a lower slot free
under one score but not the other?
\item If $A\cap B\ne\varnothing$, show that
$H(Y_{\Up A}\mid X_B)\le\chi(B)-H(X_B)$.
\end{enumerate}
\end{exercise}

Eisenstat defines perfect condensation by $\chi(B)=H(X_B)$ for every $B$.
The equivalence with the characterization above is his Theorem 5.10.
The decomposition above also explains why small score gaps imply strong
reconstruction conditions. The converse for positive-threshold almost-perfect
condensation requires additional reconstruction coverage: some active slots
may never meet that threshold. Do not identify the two approximate criteria
without checking these quantifiers.

\section{What the result says, and connections}
\label{sec:meaning}
The theorem compares two joint laws extending the same observable law. Under
the exact conditions, the two latent towers at each address determine one
another almost surely; under the quantitative hypotheses, their conditional
entropies satisfy the stated bounds. The theorem does not select the observable
family or give a learning algorithm. A perfect condensation need not exist,
as the full-support example in \cref{sec:existence} shows.

\subsection{Connections to established questions}
In statistics, a statistic is sufficient for a parameter when it
retains the data's information relevant to that parameter. Predictive
sufficiency instead concerns retaining information about a target to be
predicted. Condensation studies the organization of information across an
indexed family of observations, with explicit contribution addresses and a
correspondence result. It shares questions about information retention with
sufficiency, but these are not identical definitions or claims.

Graphical models supply the language of conditional independence and
factorization used by the ordered Markov condition. The graph here is fixed by
inclusion of addresses; its edges do not by themselves have a causal meaning.
Common-information questions also ask what two variables share. Our
correspondence argument applies across a family of overlapping address
collections, using the same entropy calculus rather than a new information
measure. Computational mechanics studies predictive states obtained from
histories with the same conditional future law. Condensation's formulation is
static: time, learning, and a prescribed past/future split are not required.
None of these comparisons implies that earlier frameworks cannot study
structured or modular organizations of information.

\subsection{Limits of information-theoretic correspondence}
If $V=f(U)$ for a bijection $f$, then
$H(U\mid V)=H(V\mid U)=0$. This remains true when describing or discovering
$f$ is expensive. The theorem does not provide an efficient translation
algorithm or account for the description length of the probability law.

Small conditional entropy is an average guarantee. If $E$ is Bernoulli$(p)$,
$D$ is an independent fair bit, $U=(E,ED)$, and $V=E$, then $H(U\mid V)=p$.
On the rare event $E=1$, the missing bit is still completely unknown.
Conversely, a rare event with an enormous number of outcomes can dominate an
entropy. A safety application needs to say which failures matter, not just
report an average information budget.

\subsection{Discussion prompts}
\begin{enumerate}
\item Choose a comparison from
\href{/interpretability/condensation-context/\#3-connections-to-prior-work}{Connections
to prior work} in the context reading. State one mathematical similarity and
one difference in what is being characterized. If it relates to one of the
condensation conditions, say how.
\item Corresponding latent collections in two models determine one another with
zero conditional entropy. What does this establish, and what else might be
needed to find and use the translation between them?
\item Where might a correspondence like this matter for AI safety? Trace one
chain from the theorem to something you would want, and say where it breaks:
which steps does the theorem not establish?
\href{/interpretability/condensation-context/\#5-the-connection-to-ai-safety}{The
connection to AI safety} sets out one such chain.
\end{enumerate}

\section{Further reading}
\label{sec:reading}
The required mathematics is contained in this sheet and the lectures.
\begin{itemize}
\item Sam Eisenstat, \emph{Condensation: A Theory of Concepts} (2025).
Definitions and exact characterization: Sections 4--5. Correspondence:
Theorems 5.15 and 6.8; two-witness inequality: Lemma 6.4.
\href{https://openreview.net/forum?id=HwKFJ3odui}{Paper and conference record}.
\item Satya Benson, \emph{Relating Almost Perfect Condensation Conditions and
the Conditioned Score}. The score/condition bridge and its reconstruction
coverage qualification.
\href{https://satyabenson.com/writing/condensation-conditions-conditioned-score}{Exposition}.
\item For the wider statistical context: Cover and Thomas,
\emph{Elements of Information Theory}; Lauritzen, \emph{Graphical Models};
Shalizi and Crutchfield, \emph{Computational Mechanics: Pattern and Prediction,
Structure and Simplicity} (2001).
\end{itemize}
\end{document}
